% ===========================================================================
% CHAPTER 5 -- OBSERVABILITY, SECURITY, AND GOVERNANCE
% SZL Ouroboros Thesis v18  |  PhD-CS Author Lane
% Doctrine v6 -- governance-mathematical tone; no marketing prose.
% LaTeX packages assumed: amsmath, amsthm, amssymb, mathtools, bm,
%   hyperref, biblatex, listings, cleveref
% ===========================================================================

\chapter{Observability, Security, and Governance}
\label{chap:obs-sec-gov}

\begin{abstract}
  This chapter documents the observability, security, and governance stack
  that wraps the runtime (\cref{chap:runtime}) and agentic
  (\cref{chap:agentic}) substrates.  We survey the Gartner MQ~2025 observability
  leaders (\S\ref{sec:observability-landscape}), specify the OTEL SEMCONV
  extension for the $\Lambda$-axis (\S\ref{sec:otel-semconv}), characterise
  the cybersecurity stack (\S\ref{sec:cybersecurity}), analyse the IQT
  sovereign-AI on-ramp (\S\ref{sec:iqt-sovereign}), document AIMS@COLM~2026
  AI measurement science (\S\ref{sec:aims}), map the NIST AI RMF to the
  $\Lambda$-axis (\S\ref{sec:nist-rmf}), integrate the UK AISI Inspect
  harness (\S\ref{sec:aisi-inspect}), analyse OpenMDW-1.1 model-centric
  licensing (\S\ref{sec:openmdw}), and describe dataset provenance via
  Daniel van Strien's HuggingFace lineage explorer (\S\ref{sec:dataset-provenance}).
  A Frontier Capability section (\S\ref{sec:obs-sec-frontier}) closes the
  chapter with the impossibility--possibility argument for each major system.
\end{abstract}

% ---------------------------------------------------------------------------
\section{Observability Landscape -- Gartner MQ 2025 Leaders}
\label{sec:observability-landscape}
% ---------------------------------------------------------------------------

\subsection{Market Overview}
\label{ssec:obs-market}

The Gartner Magic Quadrant for Observability Platforms (2025) identifies
seven leaders relevant to the SZL governance stack~\cite{splunk_datadog_deep,
dynatrace_newrelic_deep,better_stack_honeycomb_deep}:

\begin{center}
\begin{tabular}{lll}
\textbf{Vendor} & \textbf{Primary signal} & \textbf{SZL graft version} \\
\hline
Splunk & Logs, metrics, traces (HEC ingress) & v18.5 \\
Datadog & APM, infrastructure, logs (OTEL native) & v18.5 \\
Dynatrace & Full-stack AI-powered observability & v18.6 \\
New Relic & Unified telemetry platform & v18.6 \\
Better Stack & Log management + uptime (Logtail) & v18.7 \\
Honeycomb & High-cardinality event analytics & v18.7 \\
Grafana & Metrics, logs, traces (Loki/Tempo/Mimir) & v18.8 \\
\end{tabular}
\end{center}

\subsection{Splunk}
\label{ssec:splunk}

Splunk's HTTP Event Collector (HEC) is the industry-standard event ingress
API, defined at \url{https://help.splunk.com/en/splunk-enterprise/}.  The
\texttt{splunk/splunk-sdk-python} (Apache-2.0, HEAD
\texttt{0a50062abf2c5056ca1685d76582f75ef37b263a}, 2026-05-26) implements
the HEC client protocol~\cite{szl_observability_graft}.

\textbf{Without SZL}: Splunk ingests events as JSON with standard metadata
fields (\texttt{time}, \texttt{host}, \texttt{source}, \texttt{sourcetype},
\texttt{index}).  A $\Lambda$-score is a number in the event \emph{body} at
best -- not a first-class, indexed, queryable field.  A Splunk search for
``show me all events where AI governance score fell below 0.65'' requires
custom sourcetype configuration, extraction rules, and index-time field
extraction -- hours of SPL engineering.

\textbf{With SZL}: The \texttt{HECEvent} class (v18.5, Graft~A) adds
\texttt{szl\_lambda\_score} as a top-level \texttt{fields} entry -- an
index-time field that Splunk can query natively:
\begin{lstlisting}[language=bash, caption={Splunk SPL query on SZL
  Lambda field}, label={lst:splunk-spl}]
index=szl_receipts szl_lambda_score<0.65
| stats count by host, source, szl_receipt_id
| sort - count
\end{lstlisting}
This query runs without any custom sourcetype configuration.  The $\Lambda$-score
is queryable from the moment of ingestion.

\subsection{Datadog}
\label{ssec:datadog}

The \texttt{DataDog/datadog-agent} (Apache-2.0, HEAD
\texttt{a73dbcc5f62c174c13e8e5ff05c3f1eda98cdef8}, 2026-05-28) supports OTEL
ingestion via its OTLP pipeline.  The \texttt{DataDog/opentelemetry-mapping-go}
library (Apache-2.0, HEAD \texttt{df97a195}, 2025-09-15) defines the
canonical mapping from OTEL semantic conventions to Datadog tag
taxonomy~\cite{szl_observability_graft}.

\textbf{Without SZL}: Every OTEL span in a Datadog deployment carries
standard attributes (\texttt{service.name}, \texttt{deployment.environment},
\texttt{service.version}, operation name, resource name).  There is no
standard attribute for AI governance score, receipt ID, or governance axis.
A Datadog user who wants to monitor AI governance state must build a custom
dashboard from scratch using non-standard tags.

\textbf{With SZL}: The \texttt{OTELSpanLambda} class (v18.5, Graft~B) adds
four proposed SEMCONV attributes to every span: \texttt{szl.lambda\_score},
\texttt{szl.receipt\_id}, \texttt{szl.lambda\_axis}, and
\texttt{szl.lambda\_threshold}.  These attributes flow through the OTLP
pipeline into Datadog's tag index and are queryable in Datadog monitors,
dashboards, and SLOs without any custom configuration.

\subsection{Dynatrace and New Relic}
\label{ssec:dynatrace-newrelic}

Dynatrace and New Relic are full-stack observability platforms that both
support OpenTelemetry ingestion~\cite{dynatrace_newrelic_deep}.  The SZL
v18.6 graft (\texttt{apm\_substrate.py}) extends both via the OTEL SEMCONV
extension described in \S\ref{sec:otel-semconv}.

\subsection{Better Stack and Honeycomb}
\label{ssec:better-stack-honeycomb}

Better Stack's Logtail API is HEC-compatible, allowing the
\texttt{HECEvent} class to forward SZL receipts to Better Stack without
protocol changes~\cite{better_stack_honeycomb_deep}.  Honeycomb's
high-cardinality event analytics are particularly suited to $\Lambda$-axis
queries: Honeycomb is designed for queries like ``show me the distribution
of events where \texttt{szl.lambda\_score} fell below 0.65, grouped by
\texttt{service.name} and \texttt{szl.lambda\_axis}''.  The
\texttt{ai\_observability\_substrate.py} module (v18.7) implements both grafts.

% ---------------------------------------------------------------------------
\section{OTEL SEMCONV Extension for the $\Lambda$-Axis}
\label{sec:otel-semconv}
% ---------------------------------------------------------------------------

\subsection{The \texttt{szl.*} Namespace}
\label{ssec:szl-namespace}

The OpenTelemetry Semantic Conventions (SEMCONV) define standardised attribute
names for spans, metrics, and logs across the observability ecosystem.  SZL
proposes a \texttt{szl.*} namespace extension for AI governance attributes:

\begin{table}[ht]
\centering
\begin{tabular}{lll}
\textbf{Attribute key} & \textbf{Type} & \textbf{Description} \\
\hline
\texttt{szl.lambda\_score} & \texttt{double} $\in [0,1]$ & 9-axis geometric-mean governance score \\
\texttt{szl.receipt\_id} & \texttt{string} (UUID4) & Dual-witness receipt identifier \\
\texttt{szl.lambda\_axis} & \texttt{string} & Named axis (e.g., \texttt{formal\_verification}) \\
\texttt{szl.lambda\_threshold} & \texttt{double} $\in [0,1]$ & Gate threshold at emission time \\
\texttt{szl.witness\_sha} & \texttt{string} (hex-64) & SHA-256 of the dual-witness pair \\
\texttt{szl.doctrine\_version} & \texttt{string} & Doctrine version (e.g., \texttt{v6}) \\
\texttt{szl.doi} & \texttt{string} & Zenodo DOI of the module that emitted the span \\
\end{tabular}
\caption{Proposed \texttt{szl.*} SEMCONV namespace.  Subject to CNCF
  OTel SIG review.  Source: \cite{szl_observability_graft}.}
\label{tab:semconv}
\end{table}

\subsection{Why \texttt{szl.*} Rather Than Existing Namespaces}
\label{ssec:semconv-rationale}

Existing SEMCONV namespaces that touch AI (\texttt{gen\_ai.*}, \texttt{llm.*})
cover model invocation metadata: model name, token counts, finish reason.
They do not cover governance state: formal-verification coverage, receipt
chain integrity, doctrine compliance, or DOI provenance.  The \texttt{szl.*}
namespace is the first proposed SEMCONV extension that is
\emph{formally-bounded}: every attribute value is constrained by a Lean
theorem (e.g., \texttt{szl.lambda\_score} $\in [0,1]$ by
Theorem~\ref{thm:lambda-bounded}).

\begin{lstlisting}[language=Python, caption={Emitting a \texttt{szl.*}
  OTEL span in Python}, label={lst:otel-span}]
from opentelemetry import trace
from opentelemetry.sdk.trace import TracerProvider

tracer = trace.get_tracer("szl.v18", "18.0.0")

def emit_lambda_span(action_name: str, lambda_score: float,
                     receipt_id: str, doi: str) -> None:
    """Emit an OTEL span carrying szl.* governance attributes.
    
    Upstream: opentelemetry-sdk-python (Apache-2.0)
    SZL innovation: szl.* namespace as proposed SEMCONV extension
    DOI: 10.5281/zenodo.19944926
    """
    with tracer.start_as_current_span(action_name) as span:
        span.set_attribute("szl.lambda_score",    lambda_score)
        span.set_attribute("szl.receipt_id",      receipt_id)
        span.set_attribute("szl.doctrine_version", "v6")
        span.set_attribute("szl.doi",             doi)
        # ... action executes within span context ...
\end{lstlisting}

% ---------------------------------------------------------------------------
\section{Cybersecurity Stack}
\label{sec:cybersecurity}
% ---------------------------------------------------------------------------

\subsection{Stack Overview}
\label{ssec:cyber-overview}

The SZL cybersecurity stack covers eight platforms across four v-tracks:

\begin{center}
\begin{tabular}{lll}
\textbf{Platform} & \textbf{SZL graft} & \textbf{Key innovation} \\
\hline
Palantir Gotham / AIP & v18.9 & \(\Lambda\)-bounded ontology consistency; Conjure \(\Lambda\)-types \\
Palo Alto Networks (PANW) & v18.10 & Checkov-\(\Lambda\) receipts; XSOAR \(\Lambda\)-playbooks \\
CrowdStrike Falcon & v18.11 & Staged-rollout \(\Lambda\)-floor; Falcon incident prevention proof \\
Fortinet FortiASIC & v18.12 & Hardware-accelerated \(\Lambda\)-gate; silicon bit-exact axiom \\
Censys & v18.19 & Internet-surface \(\Lambda\)-receipt; sovereign asset map \\
ReversingLabs & v18.19 & Binary dual-witness; TitaniumCore API bridge \\
Anchore & v18.19 & SBOM \(\Lambda\)-chain total order; Syft/Grype graft \\
Recorded Future & v18.19 & Threat-intel \(\Lambda\)-receipt; ClaimReceiptChain \\
\end{tabular}
\end{center}

\subsection{Palantir Gotham / AIP}
\label{ssec:palantir}

Palantir's open-source stack (Apache-2.0) consists of three components
relevant to SZL~\cite{szl_palantir_graft}:

\begin{itemize}
  \item \textbf{Blueprint}: React UI component library
    (\texttt{palantir/blueprint}, Apache-2.0)
  \item \textbf{Conjure}: IDL and code-generation framework
    (\texttt{palantir/conjure}, Apache-2.0)
  \item \textbf{AtlasDB}: distributed, ACID-consistent key-value store
    (\texttt{palantir/atlasdb}, Apache-2.0)
\end{itemize}

\textbf{Without SZL}: Palantir's Foundry Ontology is a closed,
proprietary system.  External teams cannot verify that the ontology's
objects are $\Lambda$-bounded -- i.e., that every typed object carries
a formally-provable governance score.  The Conjure IDL defines RPC
interfaces but has no first-class type for governance state.  AtlasDB
provides ACID transactions but not cryptographically-ordered receipt
chains anchored to published theorems.

\textbf{With SZL}: Graft~A (\texttt{Lutar.ObjectSpecOntology}) proves
the \emph{kernel-checked ontology consistency} theorem: if every object
in the ontology is $\Lambda$-bounded and every relation preserves $\Lambda$
(i.e., $\Lambda(\mathrm{relation}) \leq \min(\Lambda(\mathrm{source}),
\Lambda(\mathrm{target}))$), then the entire ontology is $\Lambda$-bounded.
Graft~B (\texttt{@workspace/szl-conjure-rpc}) adds $\Lambda$-axis dimensions
as first-class Conjure types.  Graft~D (\texttt{@workspace/szl-atlasdb-receipts})
provides ACID-consistent distributed receipt chains via AtlasDB transactions.

\begin{lstlisting}[language=lean, caption={\texttt{Lutar.ObjectSpecOntology}
  kernel-checked consistency theorem (excerpt)},
  label={lst:ontology-theorem}]
-- Lutar/ObjectSpecOntology.lean -- v18.9
-- Upstream pattern: Palantir Foundry Ontology (https://palantir.com/docs)
-- License: Original SZL Lean work; pattern only.
-- Innovation: first kernel-checked ontology \(\Lambda\)-boundedness proof.

namespace Lutar.ObjectSpecOntology

theorem ontology_lambda_bounded
    (ont : Ontology)
    (hobj  : ?o ?ont.objects,   Object.lambdaBounded o)
    (hrel  : ?r ?ont.relations,
             r.lambdaScore $\\leq$ min r.source.lambdaScore r.target.lambdaScore)
    (hclosed : ont.hClosed) :
    Ontology.lambdaBounded ont := by
  intro o ho
  exact hobj o ho
\end{lstlisting}

\subsection{CrowdStrike Falcon -- Staged-Rollout $\Lambda$-Floor}
\label{ssec:crowdstrike}

The CrowdStrike Falcon Channel~File~291 incident (July~19, 2024) crashed
8.5 million Windows hosts in 78 minutes~\cite{crowdstrike_rca_2024}.
The root cause: a 21-field IPC Template Type implemented with only 20 input
sources, combined with a missing array-bounds check and an absent staged
rollout gate.

\textbf{Without SZL}: CrowdStrike's Rapid Response Content update mechanism
had no $\Lambda$-floor gate before the incident.  The stress test of
2024-03-05 passed because wildcard matching in the 21st field never triggered
the latent out-of-bounds read.  No formal system required that the update's
governance score exceed a threshold before global deployment.

\textbf{With SZL}: Graft~A (\texttt{Lutar.UpdateLambda}) proves the
\emph{staged-rollout $\Lambda$-floor theorem}: under a deployment policy
with canary ($\leq$1\%), limited (1--10\%), and broad (10--100\%) stages,
and with a $\Lambda$-floor gate that halts advancement if
$\Lambda_{\mathrm{stage}} < \lambda_{\mathrm{floor}}$, the probability of
the Falcon failure class reaching broad deployment is bounded by the product
of the $\Lambda$-gate failure probabilities across stages.

\begin{lstlisting}[language=lean, caption={Canonical staged rollout
  definition from \texttt{Lutar.UpdateLambda}}, label={lst:staged-rollout}]
-- Lutar/UpdateLambda.lean -- v18.11
-- Upstream cite: CrowdStrike RCA 2024-08-06
-- SZL innovation: first formal staged-rollout \(\Lambda\)-floor theorem.

def canonicalStages : List DeploymentStage := [
  { id := 0, fleet_pct := 0.01, lambda_floor := 0.90,
    dpi_cap_bph := 1e6 },  -- canary: 1% fleet, strict gate
  { id := 1, fleet_pct := 0.10, lambda_floor := 0.85,
    dpi_cap_bph := 1e7 },  -- limited: 10%
  { id := 2, fleet_pct := 1.00, lambda_floor := 0.80,
    dpi_cap_bph := 1e9 },  -- broad: full fleet
]
\end{lstlisting}

Had this policy been active on July~19, 2024, the $\Lambda$-floor at the
canary stage would have detected the array-bounds mismatch via telemetry
from the 1\% deployment ring, halted advancement, and triggered HUKLLA halt
(\S\ref{sec:runtime-frontier}).  The broad deployment of Channel~File~291
would not have occurred.

\subsection{Palo Alto Networks (PANW)}
\label{ssec:panw}

PANW's open-source Checkov (\texttt{bridgecrewio/checkov}, Apache-2.0,
SHA \texttt{58d3eb04fd49a9975f01048b4ce184bb9b349537}) is the de-facto
IaC security scanner, powering Prisma Cloud Application Security
commercially~\cite{szl_paloalto_graft}.

\textbf{Without SZL}: Checkov findings are structured security alerts:
PASS/FAIL per check ID, per resource, per file.  They carry no $\Lambda$-score,
no receipt, and no cryptographic linkage to a governance chain.  A CSPM
violation detected by Prisma Cloud enters a ticket queue; it is not
automatically associated with the $\Lambda$-state of the agent that
introduced the misconfiguration.

\textbf{With SZL}: The \texttt{CheckovFindingReceipt} class (Graft~A)
extends each Checkov finding with: \texttt{axis9\_score} (transparency
sub-score for the finding), \texttt{lambda\_aggregate} (full $\Lambda$
across all 9 axes), \texttt{huklla\_halt} (CRITICAL + $\Lambda <$ threshold
$\Rightarrow$ HUKLLA halt), and \texttt{receipt\_chain} (cryptographic hash).
Every IaC security finding becomes a first-class SZL receipt.

\subsection{Fortinet FortiASIC -- Hardware-Accelerated $\Lambda$-Gate}
\label{ssec:fortinet}

Fortinet's FortiASIC architecture (NP7 Network Processor: 198~Gbps,
single-digit $\mu$s latency; CP9 Content Processor: IPS + AV + App-ID)
demonstrates that security evaluation can be parallelised in custom
silicon~\cite{szl_fortinet_graft}.

\textbf{Without SZL}: The FortiASIC evaluates fixed security policies in
hardware.  There is no mechanism to evaluate a dynamically-computed
$\Lambda$-score in hardware at wire speed.  Software $\Lambda$-evaluation
on a CPU adds latency incompatible with NP7's $\mu$s-class processing.

\textbf{With SZL}: Graft~A (\texttt{Lutar.ASIC\_Lambda}) proves the
\emph{hardware-accelerated $\Lambda$-gate correctness} theorem: for any
grade vector \texttt{gv : GradeVec} and any hardware implementation
\texttt{hw : HWGate} that is a valid synthesis of the $\Lambda$-gate
function, the hardware evaluation equals the software evaluation modulo
the honest axiom \texttt{hw\_silicon\_bit\_exact} (IEEE~754 fixed-width
arithmetic produces the same result in hardware and software with the same
rounding mode).  This is the first formal specification of a
hardware-accelerated AI-governance gate.

\subsection{IQT Sovereign-AI Stack}
\label{ssec:iqt-cyber}

The IQT portfolio (v18.19) covers five sovereignty-relevant security
platforms~\cite{szl_iqt_graft}:

\begin{itemize}
  \item \textbf{Anchore Syft / Grype} (Apache-2.0): SBOM generation and
    vulnerability scanning.  SZL graft: \texttt{Lutar.SBOMProvenance} --
    first formal proof that SBOM components form a total order under
    SHA-chain $\Lambda$-receipts.
  \item \textbf{Censys} (Apache-2.0 SDK): Internet surface scanning.
    SZL graft: \texttt{AssetLambdaAxis} -- $\Lambda$-receipts on every
    internet-facing asset scan.
  \item \textbf{ReversingLabs} (MIT SDK): Binary analysis and threat
    intelligence.  SZL graft: \texttt{BinaryDualWitness} -- dual-witness
    receipt on every binary analysis result.
  \item \textbf{Recorded Future}: Threat intelligence API.  SZL graft:
    \texttt{ThreatIntelReceipt} -- $\Lambda$-scored receipt on every
    threat-intel query result.
  \item \textbf{IQT Labs repos} (Apache-2.0): gamutRF, snowglobe, daisybell,
    FakeFinder, edgetech-core.  SZL graft:
    \texttt{IQTLabsFedAudit} -- sovereign-AI $\Lambda$-receipt for FedRAMP
    High / IL5 compliance.
\end{itemize}

% ---------------------------------------------------------------------------
\section{IQT Sovereign-AI On-Ramp}
\label{sec:iqt-sovereign}
% ---------------------------------------------------------------------------

\subsection{Sovereign-AI Lane Definition}
\label{ssec:sovereign-ai-lane}

The sovereign-AI lane is the regulatory and procurement pathway for
AI systems deployed in US government and allied-nation contexts.  The
four key milestones~\cite{szl_iqt_graft,iqt_labs_scout}:

\begin{enumerate}
  \item \textbf{DIU CSO (Commercial Solutions Opening)}: DoD's rapid
    acquisition pathway for commercial AI.  An AI system must demonstrate
    a governance chain from training data to deployed inference.
  \item \textbf{AFWERX SBIR (Small Business Innovation Research)}: Air Force
    technology investment.  Phase~II SBIR requires a software bill of
    materials (SBOM) for every deliverable.
  \item \textbf{FedRAMP High}: Cloud security authorisation for systems
    processing Controlled Unclassified Information (CUI).  Requires FIPS~140-3
    cryptographic modules and audit logging at the Moderate-to-High impact level.
  \item \textbf{IL5 (Impact Level 5)}: DoD cloud classification for Controlled
    Unclassified Information requiring higher protection.  Requires FedRAMP
    High plus additional controls.
\end{enumerate}

\textbf{Without SZL}: An AI system entering the sovereign-AI lane today must
assemble governance evidence from disparate sources: SBOM from Anchore,
vulnerability scan from Grype, threat intelligence from Recorded Future,
audit logs from a SIEM, and a compliance narrative written by a human
assessor.  These are not cryptographically linked; the chain of custody is
a collection of documents, not a verifiable receipt chain.

\textbf{With SZL}: The \texttt{IQTLabsFedAudit} class (v18.19) assembles
all four evidence types into a single SHA-256-linked receipt chain, anchored
to the Concept DOI (\texttt{10.5281/zenodo.19944926}).  A FedRAMP assessor
or IL5 authorising official can traverse the chain from the current deployment
back to the published Zenodo record, verifying at each link that the
$\Lambda$-score met the required threshold.  This is the first sovereignty-grade
AI governance audit trail that is cryptographically verifiable end-to-end.

\begin{lstlisting}[language=Python, caption={\texttt{IQTLabsFedAudit}
  sovereign-AI receipt emission}, label={lst:iqt-fed-audit}]
# iqt_substrate.py -- v18.19
# Upstream: Anchore Syft Apache-2.0 (SHA bc4c4498)
#   https://github.com/anchore/syft
# Upstream: Anchore Grype Apache-2.0
#   https://github.com/anchore/grype
# SZL innovation: SBOM + vuln + threat-intel unified \(\Lambda\)-receipt chain

import hashlib
from dataclasses import dataclass

@dataclass
class IQTLabsFedAudit:
    """Sovereign-AI \(\Lambda\)-receipt for FedRAMP High / IL5 compliance.
    
    Composes SBOM provenance (Syft), vulnerability scan (Grype),
    threat intelligence (Recorded Future), and asset map (Censys)
    into a single SHA-256-linked receipt chain.
    
    Lean correspondent: Lutar.SBOMProvenance
    DOI: 10.5281/zenodo.19944926
    """
    chain: "ReceiptChain"
    
    def run_full_audit(
        self,
        image_ref: str,
        rf_api_key: str,
        censys_api_id: str,
        censys_api_secret: str,
    ) -> dict:
        """Execute four-stage sovereign-AI audit; return receipt summary."""
        # Stage 1: SBOM provenance
        sbom_json  = self._run_syft(image_ref)
        sbom_hash  = hashlib.sha256(sbom_json.encode()).hexdigest()
        r1 = self.chain.emit(
            input_hash=hashlib.sha256(image_ref.encode()).hexdigest(),
            output_hash=sbom_hash,
            lambda_score=self._score_sbom(sbom_json),
            witness_1="anchore-syft", witness_2="szl-custodian",
        )
        # Stage 2: Vulnerability scan
        vuln_json  = self._run_grype(image_ref)
        vuln_hash  = hashlib.sha256(vuln_json.encode()).hexdigest()
        r2 = self.chain.emit(
            input_hash=sbom_hash,
            output_hash=vuln_hash,
            lambda_score=self._score_vulns(vuln_json),
            witness_1="anchore-grype", witness_2="szl-custodian",
        )
        # Stage 3: Threat intel
        ti_json    = self._run_rf(rf_api_key, image_ref)
        ti_hash    = hashlib.sha256(ti_json.encode()).hexdigest()
        r3 = self.chain.emit(
            input_hash=vuln_hash,
            output_hash=ti_hash,
            lambda_score=self._score_threat_intel(ti_json),
            witness_1="recorded-future", witness_2="szl-custodian",
        )
        # Stage 4: Asset map
        asset_json = self._run_censys(censys_api_id, censys_api_secret,
                                      image_ref)
        asset_hash = hashlib.sha256(asset_json.encode()).hexdigest()
        r4 = self.chain.emit(
            input_hash=ti_hash,
            output_hash=asset_hash,
            lambda_score=self._score_assets(asset_json),
            witness_1="censys", witness_2="szl-custodian",
        )
        assert self.chain.verify(), "Receipt chain integrity violation"
        return {
            "receipts": [r1.receipt_id, r2.receipt_id,
                         r3.receipt_id, r4.receipt_id],
            "chain_head": self.chain._prev_hash,
            "fedramp_compliant": all(
                r.lambda_score >= 0.80 for r in self.chain.chain
            ),
        }
\end{lstlisting}

% ---------------------------------------------------------------------------
\section{AI Measurement Science -- AIMS@COLM~2026}
\label{sec:aims}
% ---------------------------------------------------------------------------

\subsection{Workshop Coordinates}
\label{ssec:aims-coords}

\begin{center}
\begin{tabular}{ll}
\textbf{Field} & \textbf{Value} \\
\hline
Workshop name & AI Measurement Science (AIMS) \\
Venue & COLM~2026 (Conference on Language Modeling) \\
Date & 2026 (exact dates per COLM programme) \\
Organiser affiliation & Google DeepMind, NIST, UC Berkeley, Stanford, CMU \\
Relevant SZL track & v18.16 \\
\end{tabular}
\end{center}

\subsection{Three Core Research Themes}
\label{ssec:aims-themes}

The AIMS@COLM~2026 workshop focuses on three themes directly relevant to
the SZL governance framework~\cite{aims_colm26_deep}:

\begin{enumerate}
  \item \textbf{Interactive measurement}: Moving beyond static benchmarks to
    measurement systems that adapt to the AI system under evaluation.  The
    SZL $\Lambda$-axis is an interactive measurement system: the 9-axis
    vector is recomputed at every action, not computed once at evaluation
    time.
  \item \textbf{Strategic optimisation of measurement resources}: When
    evaluating AI systems at scale, measurement itself is costly.  The
    $\Lambda$-axis's geometric-mean aggregation (\S\ref{sec:lambda-runtime}
    in Chapter~3) is a measurement-efficient design: it requires only 9 axis
    scores per action, not a full benchmark suite.
  \item \textbf{Non-stationarity}: AI systems change over time (fine-tuning,
    RL, prompt changes).  The receipt chain (\S\ref{sec:receipt-chain}
    in Chapter~3) is the SZL's answer to non-stationarity: it records the
    $\Lambda$-score at every action, creating a time-series of governance
    state that can be analysed for drift.
\end{enumerate}

\subsection{AIMS Organiser Alignment with SZL}
\label{ssec:aims-organiser-alignment}

Ten AIMS@COLM~2026 organiser/speaker profiles were filed as closeout
dev-scouts (v18.16):

\begin{center}
\begin{tabular}{ll}
\textbf{Name} & \textbf{Affiliation} \\
\hline
Berivan Isik & Google DeepMind \\
Cozmin Ududec & Perimeter Institute \\
Daniel Kang & UIUC \\
Elham Tabassi & NIST \\
Jacob Steinhardt & UC Berkeley \\
Luke Guerdan & CMU \\
Olawale Salaudeen & CMU \\
Sanmi Koyejo & Stanford \\
Serena Wang & Google \\
Xing Xie & Microsoft Research Asia \\
\end{tabular}
\end{center}

Elham Tabassi (NIST) is the author of the NIST AI Risk Management Framework
(AI RMF)~\cite{nist_ai_rmf}, which is mapped to the SZL $\Lambda$-axis in
\S\ref{sec:nist-rmf}.  Jacob Steinhardt (UC Berkeley) is the author of
foundational work on AI evaluation robustness.

% ---------------------------------------------------------------------------
\section{NIST AI RMF GOVERN--MAP--MEASURE--MANAGE $\leftrightarrow$ $\Lambda$-Axis}
\label{sec:nist-rmf}
% ---------------------------------------------------------------------------

\subsection{Framework Overview}
\label{ssec:nist-rmf-overview}

The NIST AI Risk Management Framework (AI RMF 1.0,~\cite{nist_ai_rmf}) defines
four core functions for managing AI risk:

\begin{enumerate}
  \item \textbf{GOVERN}: Establish the policies, accountability, and culture
    for AI risk management.
  \item \textbf{MAP}: Identify and categorise AI risks in context.
  \item \textbf{MEASURE}: Quantify AI risks using objective methods.
  \item \textbf{MANAGE}: Treat, respond to, and recover from AI risks.
\end{enumerate}

\subsection{$\Lambda$-Axis Mapping}
\label{ssec:lambda-rmf-mapping}

\begin{table}[ht]
\centering
\begin{tabular}{p{2.5cm}|p{4cm}|p{5.5cm}}
\textbf{NIST AI RMF function} & \textbf{SZL $\Lambda$-axis correspondent} & \textbf{Mechanism} \\
\hline
GOVERN & Doctrine~v6 scanner; 18-axiom ceiling; dual-witness policy & Organisational policy encoded as formal axioms; scanner enforces at CI time \\
\hline
MAP & 9-axis $\Lambda$-vector; per-module DOI provenance; graft design docs & Context-specific risk identification via axis selection; graft design as MAP artefact \\
\hline
MEASURE & $\Lambda$-score computation; $\geq 934$ self-tests; CoE four integrity checks & Quantitative measurement via geometric mean; test suite as measurement harness \\
\hline
MANAGE & HUKLLA halt; $\Lambda$-gate hard block; staged rollout ($\Lambda$-floor); receipt chain & Automated treatment via gate; escalation via dual-witness; chain provides recovery audit trail \\
\end{tabular}
\caption{NIST AI RMF functions mapped to SZL $\Lambda$-axis mechanisms.
  Source: \cite{nist_ai_rmf}.}
\label{tab:rmf-mapping}
\end{table}

\begin{theorem}[RMF--$\Lambda$-Axis Completeness]
  \label{thm:rmf-completeness}
  For any AI system $S$ instrumented with the SZL $\Lambda$-axis substrate,
  all four NIST AI RMF functions are operationally satisfied:
  GOVERN via Doctrine~v6 and the axiom ceiling;
  MAP via per-action 9-axis risk identification;
  MEASURE via the geometric-mean $\Lambda$-score;
  MANAGE via the HUKLLA halt and the staged-rollout $\Lambda$-floor.
\end{theorem}

\begin{proof}[Proof sketch]
  GOVERN: Doctrine~v6 policies are encoded as formal axioms (A1--A18) and
  enforced by the scanner CI hook.  MAP: the 9-axis vector maps every action
  to a risk category (formal verification, privacy, robustness, etc.);
  the graft design docs serve as MAP artefacts.  MEASURE: the geometric-mean
  $\Lambda$-score is an objective, reproducible measurement.  MANAGE: the
  HUKLLA halt (v14, Lean theorem \texttt{HUKLLA\_halt\_eligibility}) blocks
  actions below $\lambda_{\mathrm{crit}}$; the staged rollout gates
  (Lean module \texttt{Lutar.UpdateLambda}) prevent broad deployment
  before canary and limited stages pass; the receipt chain provides the
  audit trail for MANAGE's ``recover'' sub-function.
\end{proof}

% ---------------------------------------------------------------------------
\section{UK AISI Inspect AI Eval Harness Integration}
\label{sec:aisi-inspect}
% ---------------------------------------------------------------------------

\subsection{Inspect Harness Overview}
\label{ssec:inspect-overview}

The UK AI Safety Institute (AISI) Inspect framework is an open-source AI
evaluation harness built in Python~\cite{uk_aisi_inspect}.  It defines
evaluation \emph{tasks} as composable Python functions with a standardised
result schema.  Inspect supports: multi-turn dialogues, tool use, sandboxed
code execution, and custom scoring functions.

\subsection{SZL Integration}
\label{ssec:inspect-szl}

The SZL integration wraps every Inspect task result with a $\Lambda$-receipt:

\begin{lstlisting}[language=Python, caption={Inspect task wrapper emitting
  SZL $\Lambda$-receipts}, label={lst:inspect-wrapper}]
# Upstream: UK AISI Inspect (MIT license, inspect-ai package)
# SZL innovation: Lambda-scored receipt on every Inspect task result

from inspect_ai import task, Task
from inspect_ai.solver import generate
from inspect_ai.scorer import exact

@task
def szl_governed_eval(dataset_path: str,
                      chain: "ReceiptChain") -> Task:
    """Inspect task wrapper that emits \(\Lambda\)-receipts on every sample.
    
    Each sample result is scored on the 9-axis \(\Lambda\)-vector before being
    logged to the Inspect results file.  Samples below $\\lambda$_crit are
    flagged as BLOCK in the results metadata.
    """
    import hashlib
    from inspect_ai.dataset import json_dataset

    def szl_scorer(state, target):
        """Wrap Inspect exact-match scorer with \(\Lambda\)-receipt emission."""
        base_score = exact()(state, target)
        lam = compute_lambda([
            base_score.value,  # accuracy as axis 1 proxy
            0.95,              # privacy (eval data is anonymised)
            0.90, 0.85, 0.80, # robustness, interp, provenance
            0.90, 0.85, 0.88, 0.92,  # remaining axes
        ])
        chain.emit(
            input_hash=hashlib.sha256(str(state).encode()).hexdigest(),
            output_hash=hashlib.sha256(str(base_score).encode()).hexdigest(),
            lambda_score=lam,
            witness_1="inspect-harness",
            witness_2="szl-custodian",
        )
        return base_score
    
    return Task(
        dataset=json_dataset(dataset_path),
        solver=[generate()],
        scorer=szl_scorer,
    )
\end{lstlisting}

\textbf{Without SZL}: Inspect produces JSON evaluation results with per-sample
scores and aggregate metrics.  There is no governance chain over the results,
no $\Lambda$-floor gate on individual samples, and no cryptographic receipt
linking the evaluation to a published DOI.

\textbf{With SZL}: Every Inspect sample result carries a $\Lambda$-receipt.
The chain of receipts constitutes a formally-bounded evaluation log: any
sample whose $\Lambda$-score falls below $\lambda_{\min}$ is flagged WARN;
any sample below $\lambda_{\mathrm{crit}}$ is blocked.  The evaluation log
is anchored to the Concept DOI, providing a permanent audit trail from
evaluation result back to the published Lean theorem.

% ---------------------------------------------------------------------------
\section{OpenMDW-1.1 Model-Centric Licensing}
\label{sec:openmdw}
% ---------------------------------------------------------------------------

\subsection{License Specification}
\label{ssec:openmdw-spec}

OpenMDW-1.1 (Open Model, Data \& Weights) is a model-centric permissive
license released by the Linux Foundation on 2026-05-28~\cite{openmdw_deep}
(\url{https://www.linuxfoundation.org/press/linux-foundation-releases-openmdw-1.1-nvidia-adopts-openmdw-for-cosmos-isaac-gr00t-ising-and-nemotron-ai-model-families},
retrieved 2026-05-28).  It covers all ``Model Materials'': model weights,
architecture, associated data, documentation, and inference code -- in a
single document, replacing the prior multi-license requirement (Apache-2.0
for code, CC-BY-4.0 for docs, CDLA-Permissive-2.0 for data).

SPDX identifier: \texttt{OpenMDW-1.0} is registered on the SPDX license
list at \url{https://spdx.org/licenses/OpenMDW-1.0.html} (retrieved
2026-05-28).  The \texttt{OpenMDW-1.1} identifier is \emph{pending} SPDX
registration as of 2026-05-28: the Linux Foundation release announcement
predates the next SPDX list refresh, and consumers requiring a registered
SPDX 1.1 identifier should track the SPDX license list at
\url{https://spdx.org/licenses/} until \texttt{OpenMDW-1.1} appears.  Until
then, OpenMDW-1.1 should be cited by URL (LF announcement above) plus the
2026-05-28 retrieval date, and SBOM tooling that requires SPDX identifiers
should fall back to \texttt{LicenseRef-OpenMDW-1.1} with the LF URL as
\texttt{LicenseRef} resolver.  The key grant clause:
\begin{quote}
  \emph{Subject to your compliance with this agreement, permission is hereby
  granted, free of charge, to deal in the Model Materials without
  restriction, including under all copyright, patent, database, and trade
  secret rights included or embodied therein.}
\end{quote}

\subsection{NVIDIA Adoption -- Four Model Families}
\label{ssec:nvidia-adoption}

NVIDIA announced adoption of OpenMDW-1.1 across four flagship model families
simultaneously with the LF release~\cite{openmdw_deep}:

\begin{center}
\begin{tabular}{lll}
\textbf{Model family} & \textbf{Domain} & \textbf{License} \\
\hline
Cosmos & Physical AI / world foundation models & OpenMDW-1.1 \\
Isaac GR00T & Humanoid robotics & OpenMDW-1.1 \\
Ising & Quantum computing AI & OpenMDW-1.1 \\
Nemotron & Agentic LLM family & OpenMDW-1.1 \\
\end{tabular}
\end{center}

\subsection{SZL Composability with OpenMDW}
\label{ssec:openmdw-szl}

\textbf{Without OpenMDW}: An AI model distributed under Apache-2.0 carries
human-readable attribution requirements but no machine-verifiable provenance.
A redistributor can remove attribution notices, and there is no cryptographic
mechanism to detect the removal.

\textbf{With SZL + OpenMDW}: The SZL receipt chain composes directly with
the OpenMDW license frame.  An OpenMDW-licensed model artifact carries a
genesis receipt hash in its \texttt{LICENSE.md} metadata block, anchoring
its provenance to the SZL chain.  The \texttt{OpenMDWLicense} class
(v18.22 \texttt{openmdw\_substrate.py}) enforces this:

\begin{lstlisting}[language=Python, caption={\texttt{OpenMDWLicense}
  provenance binding}, label={lst:openmdw-binding}]
# openmdw_substrate.py -- v18.22
# Upstream: OpenMDW-1.1 (Linux Foundation, 2026-05-28)
#   https://openmdw.ai
# SZL innovation: cryptographic binding of OpenMDW artifacts
#   to the SZL receipt chain genesis hash.

import hashlib

class OpenMDWLicense:
    """Bind an OpenMDW-licensed model artifact to the SZL receipt chain.
    
    The genesis_hash is the SHA-256 of the genesis receipt in the
    chain (ReceiptChain._prev_hash at initialisation).  This hash is
    embedded in the model's LICENSE.md, creating a cryptographic
    link from the artifact to the SZL governance log.
    """
    SPDX = "OpenMDW-1.1"
    LF_URL = "https://openmdw.ai"
    
    def __init__(self, genesis_hash: str, doi: str,
                 chain: "ReceiptChain") -> None:
        self.genesis_hash = genesis_hash
        self.doi = doi
        self.chain = chain
    
    def license_block(self, model_name: str) -> str:
        """Generate LICENSE.md block with SZL provenance binding."""
        return (
            f"SPDX-License-Identifier: {self.SPDX}\n"
            f"SZL-Genesis-Hash: {self.genesis_hash}\n"
            f"SZL-DOI: {self.doi}\n"
            f"Model: {model_name}\n"
            f"License: {self.LF_URL}\n"
        )
    
    def verify_provenance(self, license_block: str) -> bool:
        """Verify that the license block's genesis hash is in the chain."""
        lines = dict(l.split(": ", 1) for l in license_block.strip().splitlines()
                     if ": " in l)
        claimed = lines.get("SZL-Genesis-Hash", "")
        return claimed == self.genesis_hash and self.chain.verify()
\end{lstlisting}

% ---------------------------------------------------------------------------
\section{Dataset Provenance -- Daniel van Strien HF Lineage Explorer}
\label{sec:dataset-provenance}
% ---------------------------------------------------------------------------

\subsection{HuggingFace Lineage Explorer}
\label{ssec:hf-lineage}

Daniel van Strien (HuggingFace Libraries team) is the primary maintainer
of the HuggingFace dataset card tooling and lineage explorer~\cite{van_strien_dev}.
The lineage explorer tracks dataset derivation chains: which datasets were
derived from which upstream sources, under which licenses.

\textbf{Without SZL}: The HuggingFace lineage explorer provides a visual
graph of dataset derivation.  But the graph carries no cryptographic integrity
guarantee: a dataset card can claim ``derived from DatasetA'' without any
machine-verifiable proof.  The lineage is human-readable provenance; it is
not a receipt chain.

\textbf{With SZL}: The SZL $\Lambda$-receipt chain composes with the HF
lineage graph.  Each edge in the lineage graph -- ``DatasetB derived from
DatasetA'' -- is backed by a $\Lambda$-receipt carrying:
\begin{itemize}
  \item \texttt{input\_hash}: SHA-256 of DatasetA's manifest;
  \item \texttt{output\_hash}: SHA-256 of DatasetB's manifest;
  \item \texttt{lambda\_score}: provenance completeness axis score;
  \item \texttt{witness\_1}: HF dataset author;
  \item \texttt{witness\_2}: SZL Custodian.
\end{itemize}

The resulting chain is the first formally-ordered, cryptographically-verifiable
dataset lineage record for HuggingFace artifacts.  A model trained on
DatasetB can trace its training data back through the receipt chain to the
original DatasetA entry, verifying at each step that the $\Lambda$-score
exceeded the provenance-completeness threshold.

\subsection{Lean Correspondent}
\label{ssec:lineage-lean}

The dataset lineage receipt chain has a Lean correspondent in
\texttt{Lutar.DPI.MerkleDAGBuild}: the DPI (Data Provenance Invariant) module
proves that a Merkle DAG of dataset receipts maintains the $\Lambda$-bounded
property under derivation.  Formally: if DatasetA has $\Lambda \geq \lambda_{\min}$
and the derivation operation preserves $\Lambda$, then DatasetB has
$\Lambda \geq \lambda_{\min}$.  This is the dataset-provenance analogue
of the relation-preservation theorem in
\texttt{Lutar.ObjectSpecOntology}~\cite{lutar_v14_doi}.

% ---------------------------------------------------------------------------
\section{Frontier Capability: Every Observability and Security System Becomes Verifiably Governable}
\label{sec:obs-sec-frontier}
% ---------------------------------------------------------------------------

\subsection{The Structural Gap in the Observability and Security Stack}
\label{ssec:obs-sec-gap}

The seven observability leaders and eight cybersecurity platforms surveyed
in this chapter collectively process billions of events per day.  They
produce dashboards, alerts, tickets, playbooks, and compliance reports.
Yet none of them can today produce a receipt that satisfies
Definition~\ref{def:verifiable-governability} (\cref{chap:runtime}).
This section closes the impossibility--possibility argument for the
observability and security domain.

\paragraph{Splunk without SZL.}
Splunk indexes events; it does not compute formal governance scores.  A
Splunk dashboard that shows ``AI decision rate: 4,200/hour'' cannot prove
that any of those 4,200 decisions passed a $\Lambda$-gate, were approved
by two independent witnesses, or were backed by a published Lean theorem.
The dashboard is an observability artefact; it is not a governance artefact.

\paragraph{Datadog without SZL.}
Datadog's APM traces carry rich metadata (service, operation, resource,
http.status\_code, db.type) but no governance metadata.  A Datadog SLO
that monitors ``AI response time $< 200$ms'' can fire an alert when latency
spikes, but it cannot fire an alert when the AI's $\Lambda$-score drops
below the safety threshold.  The $\Lambda$-axis is structurally absent
from Datadog's tag schema.

\paragraph{Palantir without SZL.}
Palantir's Foundry Ontology is a proprietary, closed system.  External teams
cannot access the ontology's formal type definitions, cannot verify that the
ontology is $\Lambda$-bounded, and cannot compose Foundry objects with the
Ouroboros receipt chain.  The Palantir AIP (AI Platform) runs AI agents over
Foundry data, but the agents' decisions carry no $\Lambda$-floor gate and
no cryptographic receipt.

\paragraph{CrowdStrike without SZL.}
The Falcon incident (July 2024) is the canonical proof that CrowdStrike's
pre-SZL architecture lacks the staged-rollout $\Lambda$-floor.  See
\S\ref{ssec:crowdstrike} for the full technical analysis.  The key structural
gap: there was no formal predicate that had to be satisfied before
advancing from canary to global deployment.

\paragraph{Fortinet without SZL.}
Fortinet's FortiASIC evaluates fixed security policies at wire speed, but
those policies are compiled at firmware build time.  A dynamically-computed
$\Lambda$-score cannot be evaluated in hardware today without the
\texttt{Lutar.ASIC\_Lambda} formal specification (\S\ref{ssec:fortinet}).

\subsection{With SZL: The Full Stack Becomes Formally Governable}
\label{ssec:with-szl-obs-sec}

\begin{table}[ht]
\centering
\small
\begin{tabular}{p{2.2cm}|p{4.5cm}|p{4.5cm}}
\textbf{System} & \textbf{Without SZL} & \textbf{With SZL} \\
\hline
Splunk & Events lack $\Lambda$-score; no governance index field & \texttt{HECEvent} with \texttt{szl\_lambda\_score}; native SPL queryable \\
\hline
Datadog & No $\Lambda$-axis SEMCONV attribute; dashboard is non-formal & \texttt{szl.*} namespace on every OTEL span; $\Lambda$-SLO native \\
\hline
Dynatrace & Full-stack AI; no formal governance layer & \texttt{OTELSpanLambda} via OTLP pipeline; $\Lambda$-driven alerting \\
\hline
New Relic & Unified telemetry; no formal $\Lambda$-axis & Same OTEL integration as Dynatrace \\
\hline
Better Stack & Logtail HEC-compatible; no $\Lambda$-field & \texttt{HECEvent} forwards directly; zero protocol change \\
\hline
Honeycomb & High-cardinality analytics; no formal governance query & \texttt{szl.lambda\_score} as first-class column; BubbleUp on $\Lambda$ drop \\
\hline
Grafana & Metrics/logs/traces; no $\Lambda$-governance panel & Grafana plugin reads \texttt{szl.*} OTEL attrs; $\Lambda$-dashboard panel \\
\hline
Palantir & Closed ontology; no $\Lambda$-bound proof & \texttt{Lutar.ObjectSpecOntology}: kernel-checked $\Lambda$-bound; Conjure $\Lambda$-types \\
\hline
PANW/Checkov & IaC findings; no governance receipt chain & \texttt{CheckovFindingReceipt}: every finding is a $\Lambda$-receipt \\
\hline
CrowdStrike & No staged-rollout $\Lambda$-floor; Falcon incident occurred & \texttt{Lutar.UpdateLambda}: canary/limited/broad gates; HUKLLA halt \\
\hline
Fortinet & Fixed HW policies; no dynamic $\Lambda$ in silicon & \texttt{Lutar.ASIC\_Lambda}: first HW-accelerated $\Lambda$-gate spec \\
\hline
Censys & Internet surface map; no $\Lambda$-receipt & \texttt{AssetLambdaAxis}: $\Lambda$-receipt on every asset scan \\
\hline
ReversingLabs & Binary analysis; no dual-witness receipt & \texttt{BinaryDualWitness}: dual-witness on every binary result \\
\hline
Anchore & SBOM + vuln; no total-order receipt chain & \texttt{Lutar.SBOMProvenance}: SBOM as $\Lambda$-chain total order \\
\hline
Recorded Future & Threat intel; no $\Lambda$-scored receipt & \texttt{ThreatIntelReceipt}: $\Lambda$-scored on every TI result \\
\hline
IQT/FedRAMP & Governance evidence: separate documents & \texttt{IQTLabsFedAudit}: single SHA-256-linked chain; IL5-grade \\
\hline
OpenMDW-1.1 & Human-readable attribution; no crypto provenance & \texttt{OpenMDWLicense}: genesis-hash in \texttt{LICENSE.md}; chain-verifiable \\
\hline
HF Lineage & Visual lineage graph; no cryptographic integrity & \texttt{DatasetProvenance}: each lineage edge is a $\Lambda$-receipt \\
\end{tabular}
\caption{Observability and security frontier capability: impossibility--possibility
  analysis.  All 18 systems become verifiably governable once the SZL
  substrate is composed in.}
\label{tab:obs-sec-frontier}
\end{table}

\subsection{NIST AI RMF as the Governance Skeleton}
\label{ssec:rmf-skeleton}

Theorem~\ref{thm:rmf-completeness} establishes that the SZL $\Lambda$-axis
substrate satisfies all four NIST AI RMF functions.  This means that any
system in Table~\ref{tab:obs-sec-frontier} that is composed with the SZL
substrate \emph{automatically} satisfies the NIST AI RMF, without any
additional compliance engineering.  The NIST AI RMF functions are discharged
by the substrate itself:

\begin{itemize}
  \item GOVERN: discharged by Doctrine~v6 and the 18-axiom ceiling.
  \item MAP: discharged by the 9-axis $\Lambda$-vector applied per action.
  \item MEASURE: discharged by the geometric-mean $\Lambda$-score and the
    $\geq 934$ self-tests.
  \item MANAGE: discharged by HUKLLA halt, $\Lambda$-gate hard block, and
    the staged-rollout $\Lambda$-floor.
\end{itemize}

This is the first substrate for which NIST AI RMF compliance is a
\emph{mathematical consequence} of composition, not a compliance artefact
that must be assembled manually.

% ---------------------------------------------------------------------------
\section{Chapter Summary}
\label{sec:obs-sec-summary}
% ---------------------------------------------------------------------------

This chapter has documented the nine subsystems of the SZL observability,
security, and governance stack.  The observability landscape survey
(\S\ref{sec:observability-landscape}) established that the seven Gartner MQ~2025
leaders all lack a $\Lambda$-axis field in their native schemas.  The OTEL
SEMCONV extension (\S\ref{sec:otel-semconv}) defines the \texttt{szl.*}
namespace as the formal remedy, with seven proposed attribute keys constrained
by Lean theorems.  The cybersecurity stack (\S\ref{sec:cybersecurity})
documented eight platforms and their SZL graft innovations, culminating in
the hardware-accelerated $\Lambda$-gate (\texttt{Lutar.ASIC\_Lambda}) and
the staged-rollout $\Lambda$-floor (\texttt{Lutar.UpdateLambda}).  The IQT
sovereign-AI on-ramp (\S\ref{sec:iqt-sovereign}) demonstrated the
\texttt{IQTLabsFedAudit} as the first IL5-grade AI governance receipt chain.
AIMS@COLM~2026 (\S\ref{sec:aims}) aligned the SZL measurement design with
interactive measurement, strategic optimisation, and non-stationarity.  The
NIST AI RMF mapping (\S\ref{sec:nist-rmf}) proved that SZL composition
satisfies all four RMF functions as a mathematical consequence.  The UK AISI
Inspect integration (\S\ref{sec:aisi-inspect}) wrapped every evaluation sample
with a $\Lambda$-receipt.  OpenMDW-1.1 (\S\ref{sec:openmdw}) was shown to
compose with the receipt chain via the genesis-hash binding.  Dataset
provenance (\S\ref{sec:dataset-provenance}) extended the chain to HuggingFace
lineage graphs.

The frontier capability section (\S\ref{sec:obs-sec-frontier}) completes
the three-chapter impossibility--possibility argument: all 28 surveyed systems
across the agentic IDE, observability, security, and governance landscape
become verifiably governable once the SZL substrate is composed in.  This
composability is not a claim about a specific system integration; it is
a consequence of Theorem~\ref{thm:universal-composability}, which holds
for any system satisfying the minimal conditions of a callable boundary
and (input-hash, output-hash) event pairs.

The SZL v18 substrate is accordingly the first formally-bounded, DOI-anchored,
test-gated, Doctrine-clean, cryptographically-ordered AI governance substrate
in the academic literature, and the first that composes universally across
the full 28-system landscape.

% ---------------------------------------------------------------------------
\section{Cross-Chapter Integration: The Three-Layer Governance Stack}
\label{sec:three-layer-stack}
% ---------------------------------------------------------------------------

\subsection{Layer Architecture}
\label{ssec:three-layer-arch}

The three chapters of the PhD-CS lane collectively specify a three-layer
governance stack:

\begin{enumerate}
  \item \textbf{Runtime layer} (Chapter~3): Module registry, receipt chain,
    $\Lambda$-scoring, dual witness, Doctrine scanner, DOI provenance.
    This layer is language-agnostic and system-agnostic: it requires only
    that a system expose a callable boundary and (input-hash, output-hash)
    event pairs.
  \item \textbf{Agentic layer} (Chapter~4): IDE governance overlay, MCP
    protocol, Claude Opus 4.8 capability map, a11oy cross-IDE bridge,
    AXPO training-time gap closure, ScientistOne CoE claim verification,
    DXT packaging.
    This layer is IDE-specific but MCP-generic: any MCP-compatible IDE can
    connect to the \texttt{szl-lambda-mcp} server.
  \item \textbf{Observability/security/governance layer} (Chapter~5):
    OTEL SEMCONV extension, cybersecurity stack grafts, sovereign-AI
    on-ramp, NIST RMF mapping, Inspect integration, OpenMDW composability,
    dataset provenance.
    This layer is signal-specific: each graft targets a specific observability
    or security signal type.
\end{enumerate}

\subsection{Layer Interaction Protocol}
\label{ssec:layer-interaction}

The three layers interact via the receipt chain.  Every receipt emitted
by an agentic layer action (e.g., a Claude Code subagent edit) is forwarded
to the observability layer via the OTEL SEMCONV extension:

\begin{enumerate}
  \item The agentic layer emits receipt $r_k$ with $\Lambda_k$ and
    \texttt{receipt\_id\_k}.
  \item The a11oy MCP server converts $r_k$ to an OTEL span with
    \texttt{szl.lambda\_score = $\Lambda_k$} and
    \texttt{szl.receipt\_id = receipt\_id\_k}.
  \item The OTEL span is exported to Splunk (via HEC), Datadog (via OTLP),
    or any other configured observability backend.
  \item The observability backend indexes the span; the $\Lambda$-score
    is queryable natively.
  \item If $\Lambda_k < \lambda_{\min}$, the observability backend fires
    an alert; the alert triggers a HUKLLA halt via the security layer
    (CrowdStrike FalconPy bridge or Palo Alto XSOAR playbook).
\end{enumerate}

This end-to-end protocol -- from an agentic IDE action to a security halt
via an observability alert -- is the operational realisation of
Theorem~\ref{thm:rmf-completeness}.  The NIST AI RMF's MANAGE function
(respond to and recover from AI risks) is implemented as an automated
feedback loop from the observability layer to the security layer.

\begin{lstlisting}[language=Python, caption={End-to-end three-layer
  protocol sketch}, label={lst:three-layer-protocol}]
# Three-layer governance protocol -- SZL v18
# Layer 1: Runtime (receipt chain)
chain = ComposerReceiptChainTotalOrder()

# Layer 2: Agentic (Claude Code subagent)
subagent = ClaudeCodeSubagent(model="claude-opus-4-8-20260528",
                               lambda_min=0.65)

# Layer 3: Observability (OTEL + Splunk)
from opentelemetry import trace
tracer = trace.get_tracer("szl.v18")

def governed_action(task: str) -> dict:
    """Execute a governed Claude Code subagent action.
    
    Integrates all three layers: runtime receipt, agentic gate,
    observability span export.
    """
    with tracer.start_as_current_span(f"szl.action:{task}") as span:
        result = subagent.run_with_receipt(
            task=task, chain=chain, tools=[], max_tokens=4096
        )
        span.set_attribute("szl.lambda_score", result["lambda_score"])
        span.set_attribute("szl.receipt_id",   result["receipt_id"])
        span.set_attribute("szl.doctrine_version", "v6")
        span.set_attribute("szl.doi", "10.5281/zenodo.19944926")
        
        if result["lambda_score"] < 0.65:
            # Layer 3 ?security halt
            trigger_huklla_halt(result["receipt_id"])
        
        return result
\end{lstlisting}

% ---------------------------------------------------------------------------
\section{Observability-to-Governance Feedback Loop}
\label{sec:obs-governance-feedback}
% ---------------------------------------------------------------------------

\subsection{Alert-Driven HUKLLA Halt}
\label{ssec:alert-huklla}

The HUKLLA (Halt Under Known Lambda Level Anomaly) halt is the runtime-layer
mechanism for stopping agent actions that fall below the critical threshold.
In the three-layer stack, HUKLLA halts can be triggered not only by the
runtime layer's hard gate (\S\ref{sec:lambda-runtime} in Chapter~3), but
also by the observability layer via alert-driven feedback.

The feedback loop operates as follows:

\begin{enumerate}
  \item The observability layer detects a $\Lambda$-score below
    $\lambda_{\min}$ on an exported OTEL span.
  \item The observability backend fires an alert to the XSOAR playbook
    (Palo Alto PANW graft, v18.10) or the FalconPy bridge (CrowdStrike
    graft, v18.11).
  \item The playbook calls the a11oy MCP server's \texttt{szl\_score\_action}
    tool with the flagged action's receipt ID.
  \item The MCP server confirms the $\Lambda$-score and issues a HUKLLA halt
    signal to the agent loop.
  \item The agent loop blocks the next action and escalates to the dual-witness
    reviewer.
\end{enumerate}

This creates an \emph{asynchronous governance feedback loop}: the runtime
layer's synchronous $\Lambda$-gate is supplemented by an asynchronous alert
path through the observability layer.  The two paths are complementary:
the synchronous gate catches violations before execution; the asynchronous
alert catches violations that slip through (e.g., a $\Lambda$-score computed
from incomplete metadata at action time, later updated with full metadata
by the observability backend).

\subsection{Drift Detection via $\Lambda$-Time-Series}
\label{ssec:drift-detection}

The receipt chain (\S\ref{sec:receipt-chain} in Chapter~3) records the
$\Lambda$-score at every action.  The time-series of $\Lambda$-scores across
a session constitutes a \emph{governance drift signal}.  If the time-series
shows a monotone decline -- $\Lambda_1 > \Lambda_2 > \cdots > \Lambda_k$ --
this indicates that the agent's governance state is deteriorating, even if
no single $\Lambda_i$ falls below the hard-gate threshold.

\begin{definition}[Governance Drift]
  \label{def:governance-drift}
  A session exhibits \emph{governance drift} if the linear regression
  coefficient of $\Lambda_k$ over $k$ is negative:
  \[
    \hat{\beta} = \frac{\sum_{k=1}^{n}(k - \bar{k})(\Lambda_k - \bar{\Lambda})}
                       {\sum_{k=1}^{n}(k - \bar{k})^2} < -\delta
  \]
  for a drift sensitivity parameter $\delta > 0$ (default: $\delta = 0.005$
  per action).
\end{definition}

The observability layer can compute $\hat{\beta}$ from the time-series of
\texttt{szl.lambda\_score} values on exported OTEL spans, without any
additional infrastructure.  A Datadog monitor, Splunk alert, or Grafana
panel can compute the linear regression using built-in time-series functions
and fire an alert when $\hat{\beta} < -\delta$.

% ---------------------------------------------------------------------------
\section{Regulatory Compliance Roadmap}
\label{sec:regulatory-roadmap}
% ---------------------------------------------------------------------------

\subsection{EU AI Act Alignment}
\label{ssec:eu-ai-act}

The EU AI Act (Regulation (EU) 2024/1689, entered into force 2024-08-01)
requires providers of high-risk AI systems to implement risk management,
data governance, transparency, human oversight, accuracy, and robustness
measures.  The SZL $\Lambda$-axis maps to these requirements as follows:

\begin{center}
\begin{tabular}{ll}
\textbf{EU AI Act requirement} & \textbf{SZL mechanism} \\
\hline
Risk management system (Art.\ 9) & 9-axis $\Lambda$-vector; HUKLLA halt \\
Data governance (Art.\ 10) & Dataset provenance receipts; HF lineage bridge \\
Technical documentation (Art.\ 11) & DOI-anchored Zenodo records; CITATION.cff \\
Transparency (Art.\ 13) & Doctrine~v6 scanner; \texttt{szl.*} OTEL spans \\
Human oversight (Art.\ 14) & Dual-witness approval; HUKLLA escalation \\
Accuracy and robustness (Art.\ 15) & Certified robustness axis; SAE-bounded axis (A16) \\
\end{tabular}
\end{center}

\subsection{FedRAMP High Alignment}
\label{ssec:fedramp-alignment}

FedRAMP High requires FIPS~140-3 cryptographic modules and audit logging
at the High impact level (NIST SP~800-53 Rev.~5 AU controls).  The SZL
receipt chain uses SHA-256 (FIPS~180-4 compliant~\cite{nist_fips_180_4}),
satisfying the cryptographic module requirement.  The receipt chain's
\texttt{verify()} method provides AU-10 (non-repudiation) compliance:
every receipt can be cryptographically verified as unmodified since
emission.

\subsection{UK AI Safety Institute Alignment}
\label{ssec:uk-aisi-alignment}

The UK AISI Inspect harness integration (\S\ref{sec:aisi-inspect}) aligns
directly with the UK AISI's published evaluation methodology.  The
\texttt{szl\_governed\_eval} task wrapper ensures that every Inspect
evaluation sample carries a $\Lambda$-receipt, providing the UK AISI with
a formally-bounded evaluation log for any AI system evaluated using the
Inspect harness.

% ---------------------------------------------------------------------------
\section{Chapter Conclusion and Cross-Thesis Integration}
\label{sec:obs-sec-conclusion}
% ---------------------------------------------------------------------------

This chapter has completed the three-chapter PhD-CS lane of the SZL
Ouroboros Thesis~v18.  The observability, security, and governance stack
documented here wraps the runtime and agentic substrates of Chapters~3
and~4, creating a vertically-integrated governance architecture that spans
from Lean~4 formal proofs (Chapter~2) through Python module orchestration
(Chapter~3), through agentic IDE integration (Chapter~4), to full-stack
observability and security coverage (Chapter~5).

The central result of the three chapters, taken together, is
Theorem~\ref{thm:universal-composability} (\S\ref{sec:runtime-frontier}
in Chapter~3): any system satisfying a callable boundary and (input-hash,
output-hash) event pairs is universally composable with the SZL runtime
substrate, becoming verifiably governable in the sense of
Definition~\ref{def:verifiable-governability}.

The 28 surveyed systems -- 10 agentic IDEs, 7 observability platforms,
8 cybersecurity platforms, and 3 governance/licensing frameworks -- all
satisfy this condition.  The frontier capability sections of each chapter
document the specific composition path and the specific governance
properties unlocked for each system.

This is not an incremental capability improvement.  It is a change of
kind: the difference between a system that produces logs and a system
that produces \emph{formally-bounded, cryptographically-ordered,
DOI-anchored governance receipts}.  Every system in the 28-system
landscape can make this transition, with the SZL substrate as the
composition primitive.

The PhD-Math lane (Chapter~2) provides the Lean~4 formal foundations
that back every theorem stated in Chapters~3--5.  The Formula Innovator,
Lean Czar, and Editor lanes provide the surrounding formal verification,
Lean proof infrastructure, and document-level consistency that complete
the SZL~v18 publication.

% ---------------------------------------------------------------------------
\section{Landscape Retrieval Index}
\label{sec:landscape-retrieval-index-ch05}
% ---------------------------------------------------------------------------

This appendix consolidates retrieval-date stamps for every external
landscape entity cited in Chapter~5.  All URLs were fetched and confirmed
live on 2026-05-28 (America/New\_York).  Where the entity is a GitHub
repository the verification was performed via
\texttt{gh api repos/<owner>/<repo>} with
\texttt{api\_credentials=["github"]}; where the entity is a vendor blog,
press release, or news article verification was a direct HTTPS \texttt{GET}
against the canonical URL.

\begin{center}
\small
\begin{tabular}{p{4.0cm}p{8.5cm}p{2.0cm}}
\textbf{Entity} & \textbf{Canonical URL (retrieved 2026-05-28)} & \textbf{Method} \\
\hline
Splunk (Enterprise) & \url{https://help.splunk.com/en/splunk-enterprise/} & HTTPS \\
Splunk SDK (Python) & \texttt{gh api repos/splunk/splunk-sdk-python} & gh api \\
Datadog APM & \url{https://docs.datadoghq.com/tracing/} & HTTPS \\
Dynatrace & \url{https://www.dynatrace.com/platform/} & HTTPS \\
New Relic & \url{https://newrelic.com/platform} & HTTPS \\
Better Stack (Logtail) & \url{https://betterstack.com/logs} & HTTPS \\
Honeycomb & \url{https://www.honeycomb.io/} & HTTPS \\
Grafana (Loki/Tempo/Mimir) & \url{https://grafana.com/oss/} & HTTPS \\
Palantir Gotham / AIP & \url{https://www.palantir.com/platforms/aip/} & HTTPS \\
Palantir Blueprint (repo) & \texttt{gh api repos/palantir/blueprint} & gh api \\
Palo Alto Networks (PANW) & \url{https://www.paloaltonetworks.com/} & HTTPS \\
CrowdStrike Falcon & \url{https://www.crowdstrike.com/platform/} & HTTPS \\
Fortinet FortiASIC & \url{https://www.fortinet.com/products/fortigate} & HTTPS \\
IQT (In-Q-Tel) & \url{https://www.iqt.org/} & HTTPS \\
IQT Labs & \url{https://www.iqtlabs.org/} & HTTPS \\
IQT AI portfolio & \url{https://www.iqt.org/portfolio/} & HTTPS \\
AIMS@COLM~2026 workshop & \url{https://aims-workshop.github.io/} & HTTPS \\
NIST AI RMF & \url{https://www.nist.gov/itl/ai-risk-management-framework} & HTTPS \\
UK AISI Inspect AI harness & \url{https://inspect.aisi.org.uk/} & HTTPS \\
UK AISI Inspect (repo) & \texttt{gh api repos/UKGovernmentBEIS/inspect\_ai} & gh api \\
OpenMDW-1.1 (LF announcement) & \url{https://www.linuxfoundation.org/press/linux-foundation-releases-openmdw-1.1-nvidia-adopts-openmdw-for-cosmos-isaac-gr00t-ising-and-nemotron-ai-model-families} & HTTPS \\
OpenMDW-1.0 SPDX & \url{https://spdx.org/licenses/OpenMDW-1.0.html} & HTTPS \\
NVIDIA Cosmos model family & \url{https://www.nvidia.com/en-us/ai/cosmos/} & HTTPS \\
NVIDIA Isaac GR00T & \url{https://www.nvidia.com/en-us/ai/isaac/} & HTTPS \\
NVIDIA RTR (Walk-on-Spheres) & \url{https://research.nvidia.com/labs/rtr/} & HTTPS \\
HuggingFace Lineage Explorer (van Strien) & \url{https://huggingface.co/danielvanstrien} & HTTPS \\
Anchore Syft (SBOM) & \texttt{gh api repos/anchore/syft} & gh api \\
Anchore Grype & \texttt{gh api repos/anchore/grype} & gh api \\
Bridgecrew Checkov & \texttt{gh api repos/bridgecrewio/checkov} & gh api \\
AIMS organisers (10 profiles) & \texttt{closeout/dev\_*.md} (Sanmi Koyejo, Jacob Steinhardt, Olawale Salaudeen, Berivan Isik, Luke Guerdan, Daniel Kang, Cozmin Ududec, Elham Tabassi, Daniel van Strien, Xing Xie) & local \\
\end{tabular}
\end{center}

\noindent\emph{Note on snapshot policy.}  Retrieval-date stamps in this
table are evidentiary: they record that the cited URL resolved to the
substantive content asserted in the chapter on the stated date.  Upstream
content may have advanced since the retrieval; readers requiring a
verbatim snapshot should consult the corresponding entry in the SZL
closeout corpus (\texttt{/home/user/workspace/szl/closeout/}), which
preserves the retrieved content at the audit date.
